\topic{sutton-definitions}{Sutton Section 2.1: the definitions, and the shorthand behind them}

Reading anchor: Sutton and Biblarz, ninth edition, Section 2.1,
``Definitions'', printed pp.~26--31 / PDF pp.~50--55, Eqs.~2--1 to 2--11 and
Example~2--1 \cite{sutton}. The text reproduced below is Ernest's own
section on the same material from \path{LaTeXandpdfs/Propulsion.tex}, written
against the seventh edition \cite{GSuttonOBiblarz2001}; the equation numbering
of Section 2.1 and the placement of Example~2--1 are unchanged between the two
editions, so its citations resolve in either. (The book's authors are Sutton
and Biblarz in that order; the manuscript's ``Biblarz and Sutton'' is kept as
written.)

The point of carrying this section is the bridge it builds. The previous topic,
Section~\ref{topic:force-one-form}, is the rigorous end: force is a covector,
impulse is its increment, and Sutton's scalars are components taken under
stated assumptions. The manuscript below is the engineering end: definitions
boxed, the practical formulas recovered, the mass bookkeeping made explicit
with ``imaginary boxes''. What is added here sits between them. Each block of
the manuscript is followed by run-in paragraphs headed \textbf{Annotation},
which do one of three things: name the assumption a shortcut is silently
using, check a claim, or connect a formula to the covector picture. The
manuscript text itself is not altered except for the transcription fixes
listed at the end; every correction of substance is made in an annotation,
where the reader can see it.

\subsection{Definitions}

Section 2.1. ``Definitions'', Ch. 2 Definitions and Fundamentals of Biblarz
and Sutton (2001) \cite{GSuttonOBiblarz2001}

The \textbf{total impulse} $I_t$ is the thrust force
$F_{\text{thrust}} = F_{\text{thrust}}(t)$ integrated over the burning time
$T$, i.e.

total impulse $I_t = \int_0^T F_{\text{thrust}} dt $ \\
$t\equiv $ burning time $ \equiv t_p$

specific impulse $I_s$, total impulse per unit weight of propellant

\[
I_s = \frac{ \int_0^t F_{\text{thrust}} dt }{ g_0 \int \dot{m} dt }
 = \frac{I_t}{ m_p g_0}
\]
\[
I_s = \frac{ \dot{m} u_e t_p }{ \dot{m} g_0 t_p } = \frac{u_e}{g_0}
\]
If $F_{\text{thrust}} = \dot{m} u_e$, and constant propellant mass flow,
$u_e$.

Instead, define
\begin{definition}[total impulse]
\begin{equation}\label{Eq:TotalImpulseDefinition}
\boxed{ I_{\text{tot}} \equiv I_t = \int_0^t dt F_{\text{thrust}}(t)    }
\end{equation}
\end{definition}

\paragraph*{Annotation: what the quick lines assume.}
The two displayed lines before the boxed definition are Sutton's Eqs.~2--3
and 2--5 with the route between them compressed. Written out, the second line
uses four things at once: $\dot m$ constant, $u_e$ constant, no start or stop
transients, and $F=\dot m u_e$ with nothing else --- that is, no pressure
thrust. The first three are A2--A4 of Section~\ref{topic:force-one-form}; the
fourth is discussed under the exhaust-velocity annotation below, and it is why
Sutton's $u_e$ in this role is the \emph{effective} exhaust velocity $c$, not
the nozzle-exit velocity $v_2$. The $g$ that divides is the defined constant
$g_0=9.80665\ \mathrm{m/s^2}$, never the local gravitational acceleration ---
Sutton says so explicitly, and it is why $I_s$ ``in seconds'' is not a
duration. Three names for the burn time ($T$, $t$, $t_p$) are in use;
Sutton uses $t$ inside the integrals of Chapter~2 and $t_p$ in Chapter~4.

\paragraph*{Annotation: the box is a component of a covector.}
Equation~\ref{Eq:TotalImpulseDefinition} is a scalar. In the language of
Section~\ref{topic:force-one-form}, it is $J(\hat e)$, the impulse covector
$J=\int f\,dt$ paired with the fixed thrust axis $\hat e$, and the pairing is
what makes it a number rather than an object living in different cotangent
spaces at different instants. It equals the magnitude of the delivered
momentum change only when the thrust direction holds still
(Proposition~\ref{prop:scalar-bound}); for a gimballed or slewing burn the
boxed number is an upper bound on $\lVert\Delta p\rVert$, not the $\Delta v$
budget.

Let $m_{p, \text{exp}} = m_{p, \text{exp}}(t) = $ mass of propellant already
expelled by time $t$. $m_{p, \text{exp}}(t) \in C^{\infty}(\mathbb{R})$.

Biblarz and Sutton (2001) \cite{GSuttonOBiblarz2001} defined \emph{specific
impulse} to be the total impulse per unit weight of propellant as such:
\begin{definition}[Specific Impulse]
\begin{equation}\label{Eq:SpecificImpulseDefinition}
\boxed{ I_s := \frac{  \int_0^t dt F_{\text{thrust}}(t) }{ W_{p, \text{total expelled} }} }
\end{equation}
and so
\begin{equation}\label{Eq:SpecificImpulseIntegralForm}
I_s = \frac{\int_0^t dt F_{\text{thrust}}(t) }{ g_0 \int dt \dot{m}_{p, \text{exp}}(t)}
\end{equation}
For practical calculations,
\begin{equation}\label{Eq:SpecificImpulsePractical}
I_s = \frac{ \overline{F}_{\text{thrust} } }{ g_0 \overline{\dot{m}}_{p, \text{expelled}}  }
\end{equation}
\end{definition}
so it's more honest to say that it's the time-averaged specific impulse.

Note that $\int dt \dot{m}_{p, \text{exp}}(t) = m_{p, \text{exp}} = $ total
effective propellant mass expelled through the nozzle.

\paragraph*{Annotation: which average it is.}
Two separate points. First,
Eq.~\ref{Eq:SpecificImpulsePractical} is not an approximation to
Eq.~\ref{Eq:SpecificImpulseIntegralForm}; it is the same number, provided both
bars average over the same window $[0,t]$:
\[
 \frac{\overline F}{g_0\,\overline{\dot m}}
 =\frac{\frac1t\int_0^tF\,dt}{g_0\,\frac1t\int_0^t\dot m\,dt}
 =\frac{\int_0^tF\,dt}{g_0\int_0^t\dot m\,dt},
\]
the $1/t$ cancelling. Second, and this is where ``time-averaged'' misleads,
the number is \emph{not} the time average of the instantaneous specific
impulse $I_s(t):=F(t)/(\dot m(t)g_0)$. Substituting $F=I_s(t)\,\dot m\,g_0$
into the numerator,
\[
 \frac{\int_0^tF\,dt}{g_0\int_0^t\dot m\,dt}
 =\frac{\int_0^tI_s(t)\,\dot m(t)\,dt}{\int_0^t\dot m(t)\,dt}
 \;\ne\;\frac1t\int_0^tI_s(t)\,dt
 \quad\text{unless $\dot m$ (or $I_s$) is constant,}
\]
which is Proposition~\ref{prop:weighted-mean}: the definition is the
mass-flow-weighted mean of $I_s(t)$, the average that weights each second of
burn by how much propellant it spent. Sutton's own phrase for Eq.~2--3,
``a time-averaged specific impulse value'', is loose in exactly this way. The
honest sentence is: $I_s$ as defined is the ratio of two totals, and equals
the time average of $I_s(t)$ only for constant mass flow. For a
boost--sustain motor with $(F_1,\dot m_1)$ over $t_1$ and $(F_2,\dot m_2)$
over $t_2$, the two averages are
$\frac{I_{s,1}\dot m_1t_1+I_{s,2}\dot m_2t_2}{\dot m_1t_1+\dot m_2t_2}$ and
$\frac{I_{s,1}t_1+I_{s,2}t_2}{t_1+t_2}$, and the first is the one that
predicts $\Delta v$.

\paragraph*{Annotation: weight, and regularity.}
$W_{p,\text{total expelled}}=g_0\,m_{p,\text{exp}}$ is a mass multiplied by a
conventional constant. At sea level it is a weight; in orbit Sutton notes it
``does not represent the weight'' \cite[Section 2.1]{sutton}. Nothing in the
definition uses a gravitational field. On regularity: a real burn has an
ignition transient and a shutdown transient, and $m_{p,\text{exp}}$ is
honestly piecewise $C^1$, with $\dot m_{p,\text{exp}}\ge0$ everywhere and
$=0$ outside the burn. Assuming $C^\infty$ is a convenience (mollify the
transients), and it is harmless for every formula here, all of which are
integrals. The one place it bites is the strict inequality
$\dot m_{p,\text{exp}}>0$ written below: with $m_{p,\text{exp}}(0)=0$ and the
expelled mass nonnegative, the derivative must vanish for $t\le0$, so the
correct statement is $\ge0$ with equality before ignition and after burnout.

If we are equipped with the exterior covariant derivative
$D: \Omega^{d-1}(N; TN) \to \Omega^d(N; TN)$, then the force on a fluid over
any general manifold is
\[
F = \int_B \dot{m} \otimes \mathbf{u} + m \otimes \left( \frac{\partial \mathbf{u}}{ \partial t} + \nabla_{\mathbf{u}} \mathbf{u} \right)
\]

If $\frac{\partial \mathbf{u}}{ \partial t} = 0$,
$\nabla_{\mathbf{u}} \mathbf{u} = 0$ (no curved space),
\[
F = \int_B \dot{m} \otimes \mathbf{u}
\]
If $\mathbf{u}$ constant over $B$, call the \emph{effective exhaust velocity}
$c \equiv u$ to be
\begin{equation}\label{Eq:ConstantEffectiveExhaustVelocity}
\begin{gathered}
F = \dot{m} u  \\
u = \frac{F}{\dot{m}} = I_s g_0
\end{gathered}
\end{equation}
cf. Eq. 2-6 of Biblarz and Sutton (2001) \cite{GSuttonOBiblarz2001}

\paragraph*{Annotation: what that expression is, and what replaces it.}
The displayed $F$ is the product rule applied to ``$m\mathbf u$'' and
integrated: $\dot m\,\mathbf u$ plus $m$ times the material acceleration
$\frac{\partial\mathbf u}{\partial t}+\nabla_{\mathbf u}\mathbf u$. It is a schematic, not a
derivation, and it should be read as a signpost to the right object rather
than as a formula to evaluate. Three remarks make it precise.

First, the notation $D:\Omega^{d-1}(N;TN)\to\Omega^d(N;TN)$ names the correct
tool: a momentum flux through a surface is a vector-valued $(d-1)$-form, and
its exterior covariant derivative is the coordinate-free way to state a
momentum balance on a volume. That machinery is what
Section~\ref{topic:force-one-form} defers to a Sutton~2.2 topic; the
displayed formula does not actually use it.

Second, the reduction ``$\nabla_{\mathbf u}\mathbf u=0$ (no curved space)''
mislabels a flow assumption as a geometric one. In flat space
$\nabla_{\mathbf u}\mathbf u=(\mathbf u\cdot\nabla)\mathbf u$ is the
convective acceleration, and it vanishes only when streamlines are straight
and the speed is constant along them. That is precisely the assumption Sutton
attaches to Eq.~2--12: ``the exit gas velocity is assumed constant, uniform,
and purely axial.'' Curvature of space is not in play anywhere in this
chapter.

Third, the surviving term $\int_B\dot m\otimes\mathbf u$ is a momentum
\emph{flux through a surface}, not a derivative of a vehicle's momentum, and
Section~\ref{topic:force-one-form} already warns that the vehicle-only
$F=d(mv)/dt$ is not a correct starting point for a variable-mass body. The
manuscript's later remark that ``$\dot m_S$ is the mass per unit (surface)
area'' is the right instinct with the wrong units: the integrand is the mass
\emph{flow rate} per unit area, $\dot m_S=\rho\,\mathbf u\cdot\mathbf n$ in
$\mathrm{kg\,m^{-2}\,s^{-1}}$, and $\int_S\dot m_S\,dA=\dot m$. The honest
flat-space statement is the control-volume momentum balance, which for a fixed
control volume $B$ with outward normal $\mathbf n$, Cauchy traction
$\mathbf t=\boldsymbol\sigma\mathbf n$ and body force $\mathbf b$ reads, in
Cartesian components,
\begin{equation}\label{Eq:ControlVolumeMomentumBalance}
 \frac{d}{dt}\int_B\rho\,u^i\,dV
 +\oint_{\partial B}\rho\,u^i\,(u^jn_j)\,dA
 =\oint_{\partial B}t^i\,dA+\int_B\rho\,b^i\,dV .
\end{equation}
Both terms of the manuscript's schematic are inside the left side: by the
transport theorem and continuity, for a steady flow
$\int_B\rho\,(\mathbf u\cdot\nabla)u^i\,dV=\oint_{\partial B}\rho u^i(u^jn_j)\,dA$,
so ``$\dot m\otimes\mathbf u$'' and ``$m\otimes\nabla_{\mathbf u}\mathbf u$''
are the same quantity in two forms, and adding them would count it twice.

\paragraph*{Annotation: from the balance to the thrust equation.}
Sutton's ninth edition quotes rather than derives its thrust equation
(``a derivation can be found in earlier editions''), so here is the
control-volume derivation in full, for the test-stand case that his Eq.~2--13
describes. Let $\hat e$ be the forward unit vector. Take $B$ to be the region
bounded by a surface $\Sigma$ that hugs the outside of the vehicle and cuts
straight across the nozzle exit plane $A_2$, so the contents of $B$ are the
vehicle, the propellant aboard, and the gas in the chamber and nozzle. Assume:
steady flow in the vehicle frame, which is inertial on a test stand
(A\textsubscript{s}); exit velocity uniform, axial, of magnitude $v_2$
(A\textsubscript{u}); uniform ambient pressure $p_3$ on the outside of
$\Sigma$ and uniform gas pressure $p_2$ across $A_2$ (A\textsubscript{p});
inviscid traction on $\Sigma$, $t^i=-p\,n^i$ (A\textsubscript{v}); and no body
force, gravity being balanced separately or absent (A\textsubscript{g}). Under
A\textsubscript{s} the first term of
Eq.~\ref{Eq:ControlVolumeMomentumBalance} vanishes. Mass crosses $\Sigma$ only
through $A_2$, where $\mathbf u=-v_2\hat e$ and $\mathbf n=-\hat e$, so with
A\textsubscript{u} the flux term is
\[
 \oint_{\partial B}\rho\,\mathbf u\,(\mathbf u\cdot\mathbf n)\,dA
 =\int_{A_2}\rho_2\,(-v_2\hat e)\,(v_2)\,dA
 =-\rho_2v_2^2A_2\,\hat e=-\dot m\,v_2\,\hat e,
 \qquad \dot m=\rho_2v_2A_2 .
\]
The external forces on the contents are the stand's reaction $\mathbf R$ and
the pressure over $\Sigma$. Write the pressure as $p_3$ everywhere on $\Sigma$
except on $A_2$, where it is $p_2$, and use A\textsubscript{v}:
\[
 \oint_{\partial B}\mathbf t\,dA
 =-\oint_{\Sigma}p\,\mathbf n\,dA
 =-\underbrace{\oint_{\Sigma}p_3\,\mathbf n\,dA}_{=\,0\ \text{(closed surface)}}
  -\int_{A_2}(p_2-p_3)\,\mathbf n\,dA
 =(p_2-p_3)A_2\,\hat e .
\]
Equation~\ref{Eq:ControlVolumeMomentumBalance} then reads
$-\dot m v_2\hat e=\mathbf R+(p_2-p_3)A_2\hat e$, i.e.
\[
 \mathbf R=-\bigl[\dot m\,v_2+(p_2-p_3)A_2\bigr]\hat e ,
\]
the stand pushes backward with exactly the thrust, and by Newton's third law
the vehicle pushes forward with
\begin{equation}\label{Eq:ThrustEquation}
 F=\dot m\,v_2+(p_2-p_3)A_2 ,
\end{equation}
Sutton's Eq.~2--13. The momentum thrust is the flux term; the pressure thrust
is the one place where the closed-surface cancellation of $p_3$ fails, namely
on the cut where the pressure is $p_2$ instead. Dropping the pressure term
requires $p_2=p_3$ (optimum expansion, Chapter~3), which is the additional
assumption hidden in the manuscript's $F=\dot m u$.

\paragraph*{Annotation: why the boxed equation is nevertheless exact.}
The subtlety is that Eq.~\ref{Eq:ConstantEffectiveExhaustVelocity} is true
\emph{by construction}, with no flow assumption at all. Sutton does not
measure $c$; he defines it, Eq.~2--6, as the number $F/\dot m$, and then
observes, Eq.~2--15, that Eq.~\ref{Eq:ThrustEquation} gives it the value
\[
 c=\frac{F}{\dot m}=v_2+\frac{(p_2-p_3)A_2}{\dot m}=I_sg_0 .
\]
So $F=\dot mc$ holds whenever $\dot m$ is constant, full stop, and
A\textsubscript{u}, A\textsubscript{p} and $p_2=p_3$ are needed only to
identify the defined $c$ with the physical exit velocity $v_2$. This is the
engineering shorthand at its most honest: $c$ is the velocity that
\emph{would} produce the measured thrust if all of it were momentum thrust.
What the covector picture adds is that even this is a metric statement ---
$F/\dot m$ is a covector divided by a scalar, and calling it a velocity is
$c\,\hat e=(f/\dot m)^\sharp$, the musical isomorphism applied silently, as
Section~\ref{topic:force-one-form} spells out.

\textbf{Mass ratio} $\mathbf{MR}$ of a vehicle or \emph{particular vehicle
stage} is defined as
\[
\mathbf{MR} := \frac{m_f}{m_0}
\]
where, using Biblarz and Sutton (2001)'s notation, \\
$m_f \equiv $ final mass (after rocket operation has consumed all usable
propellant), and \\
$m_0$ (before rocket operation).

$\mathbf{MR}$ applies to a single, or multistage vehicle.
\begin{itemize}
	\item Overall mass ratio is the product of the individual vehicle stage
	mass ratios
	\item Final mass $m_f$ is mass of vehicle after rocket has ceased to
	operate when all useful propellant mass $m_p$ has been consumed and
	ejected.
	\item $m_f$ includes all components not useful propellant and may include
	guidance devices, navigation gear, \textbf{payload}
	\end{itemize}

\paragraph*{Annotation: the product of stage mass ratios.}
The first bullet is Sutton's sentence, and it is true in the sense that
matters for velocity, not literally for masses. Take $N$ stages burned in
series; stage $i$ starts at $m_{0,i}$ and ends at $m_{f,i}$, then jettisons
its inert hardware $s_i$, so $m_{0,i+1}=m_{f,i}-s_i$. Telescoping,
\[
 \frac{m_{f,N}}{m_{0,1}}
 =\prod_{i=1}^{N}\frac{m_{f,i}}{m_{0,i}}\;
  \prod_{i=1}^{N-1}\frac{m_{0,i+1}}{m_{f,i}}
 =\Bigl(\prod_{i=1}^{N}\mathbf{MR}_i\Bigr)
  \prod_{i=1}^{N-1}\Bigl(1-\frac{s_i}{m_{f,i}}\Bigr),
\]
so the literal ratio of final to initial mass equals the product of the stage
ratios only when nothing is dropped. What \emph{is} multiplicative is the
quantity the mass ratio is for: in gravity-free, drag-free flight each stage
contributes $\Delta u_i=c_i\ln(1/\mathbf{MR}_i)$ (Sutton Eq.~4--6), and with
equal $c$ the total is $c\ln\prod_i(1/\mathbf{MR}_i)$. Sutton's ``overall
mass ratio'' is the product in that sense, made precise in his Section~4.7,
and this is the first place the manuscript's imaginary boxes earn their keep:
each $\mathbf{MR}_i$ is computed in the box of stage $i$ with everything above
it counted as payload, and the boxes for successive stages are not nested
copies of one another.

\subsubsection{Mass definitions for rockets, vehicles, propellants}

Instead, let the mass of the rocket $+$ remaining propellant be
$M = M(t) \in C^{\infty}(\mathbb{R})$, such that $M \geq 0$,
$\forall \, t \in \mathbb{R}$, always.

Let $M_f \in \mathbb{R}^+$ be a constant, what Biblarz and Sutton calls the
final or inert propulsion mass, and denote as $m_f$.

Let $m_p = m_p(t) = $ propellant remaining on vehicle or stage. Thus,
\begin{equation}
\dot{m}_p \leq 0
\end{equation}
which is expected by physical considerations.

Let
\begin{equation}\label{Eq:PropellantRemainingBookkeeping}
m_p(t)  = m_p(0) - m_{p,\text{exp}}(t)
\end{equation}
where $m_{p,\text{exp}}(t) = $ total propellant that had been expelled.
Clearly
\begin{equation}
\dot{m}_{p, \text{exp}}(t) \geq 0, \, m_{p,\text{exp}}(0) = 0
\end{equation}

Thus
\begin{equation}\label{Eq:MassBookkeeping}
\begin{aligned}
& M(t) = M_f + m_p(t)  \\
& M(0) = M_f + m_p(0)
\end{aligned}
\end{equation}
and so define
\begin{definition}[Mass ratio]
\begin{equation}\label{Eq:MassRatioDefinition}
\mathbf{MR} := \frac{M_f}{M(0)}	= \frac{M_f}{ M_f + m_p(0)}
\end{equation}
	\end{definition}
and
\begin{definition}[propellant mass fraction $\zeta$]
	\begin{equation}\label{Eq:PropellantMassFractionDefinition}
	\zeta := \frac{m_p(0)}{ M(0)} = \frac{M(0) - M_f}{M(0)}
	 = \frac{m_p(0)}{M_f + m_p(0)}
	\end{equation}
	\end{definition}
Biblarz and Sutton (2001) uses a different notation for the masses.

e.g. $\zeta = \frac{m_p(0)}{ M(0)} = 0.91$ means only 9 percent of mass is
inert rocket hardware, and this small fraction contains, feeds, burns a
substantially larger mass of propellant: a high value of $\zeta$ is
desirable.

\paragraph*{Annotation: the dictionary, and one identity.}
The manuscript's time-dependent bookkeeping and Sutton's static symbols
correspond as follows, with the box always to be stated:
\[
 M(0)\leftrightarrow m_0,\qquad
 M_f\leftrightarrow m_f,\qquad
 m_p(0)\leftrightarrow m_p,\qquad
 M(t_p)=M_f .
\]
For the whole vehicle $M_f=M_r+M_{\text{pay}}$; for the propulsion system
alone $M_f=M_r$ (the manuscript's notation for Example~2--1, below). Sutton's
Eqs.~2--7 to 2--10 are, in this notation,
Eq.~\ref{Eq:MassRatioDefinition}, Eq.~\ref{Eq:PropellantMassFractionDefinition}
and the second line of Eq.~\ref{Eq:MassBookkeeping}; the one identity worth
keeping in view is that the two definitions are not independent,
\begin{equation}\label{Eq:ZetaPlusMR}
 \zeta+\mathbf{MR}
 =\frac{m_p(0)}{M(0)}+\frac{M_f}{M(0)}
 =\frac{m_p(0)+M_f}{M(0)}=1 ,
\end{equation}
so a ``mass fraction of 0.91'' and a ``mass ratio of 0.09'' are the same
statement, and $1/\mathbf{MR}=1/(1-\zeta)$ is the number that enters the
logarithm of the rocket equation. The residual-propellant convention matters
here: Sutton counts unusable residuals in $m_f$, some manufacturers count
them in $m_p$, and $\zeta$ moves accordingly.

Impulse to weight ratio of a complete propulsion system:

\[
\frac{I_t}{w_0}
\]
$I_t \equiv $ total impulse, \\
$w_0 \equiv $ initial or propellant-loaded vehicle weight.

It's reasonable to suppose that
\begin{equation}\label{Eq:ImpulseToWeight}
\begin{gathered}
\frac{I_t}{w_0} = \frac{I_t}{ M(0) g_0}
 = \frac{I_t}{  (M_f + m_p(0)) g_0}
 = \frac{I_s m_p(0) g_0}{ (M_f + m_p(0)) g_0}
 = \frac{ I_s}{ \frac{M_f}{m_p} + 1}
\end{gathered}
\end{equation}

thrust-to-weight ratio $F/w_0$

\paragraph*{Annotation: not a supposition, and a cleaner form.}
Equation~\ref{Eq:ImpulseToWeight} is Sutton's Eq.~2--11, and nothing in it
needs supposing. The only input beyond the definitions is
$I_t=I_sm_p(0)g_0$, which is Eq.~\ref{Eq:SpecificImpulseIntegralForm}
rearranged (Sutton's Eq.~2--4) and is therefore exact for any thrust
history --- the ``constant thrust, negligible transients'' Sutton attaches to
Eq.~2--11 is inherited from reading $I_s$ as Eq.~2--5 and is not needed when
$I_s$ is the integral definition. The last member also has a cleaner name:
\begin{equation}\label{Eq:ImpulseToWeightIsIsZeta}
 \frac{I_t}{w_0}
 =I_s\,\frac{m_p(0)}{M_f+m_p(0)}
 =I_s\,\zeta ,
\end{equation}
by Eq.~\ref{Eq:PropellantMassFractionDefinition}. So the impulse-to-weight
ratio is the specific impulse discounted by the propellant fraction, it has
the same conventional unit of seconds, and it satisfies $I_t/w_0\le I_s$ with
equality only for a vehicle made entirely of propellant. That is the exact
content of Sutton's remark, in Example~2--1, that a ``good design'' would
bring $I_t/w_0$ close to $I_s$: it means $\zeta$ close to one.

For the thrust-to-weight ratio, with $F$ constant and gravity and drag
neglected, $F/w_0=a_0/g_0$ is the initial acceleration in units of $g_0$, and
because the thrust is unchanged while the mass falls,
\begin{equation}\label{Eq:BurnoutAcceleration}
 a_f=\frac{F}{M_f}=\frac{F}{M(0)}\,\frac{M(0)}{M_f}
 =\frac{a_0}{\mathbf{MR}} ,
\end{equation}
the burnout acceleration is the initial one divided by the mass ratio of
whichever box actually moves together.

Example 2-1 of Biblarz and Sutton (2001) \cite{GSuttonOBiblarz2001} has some
subtleties about the mass definitions that wasn't covered clearly (at least
for me) by the discussion.

Let $M_r$ is the mass of the propulsion system, $r$ standing for ``rocket.''
If there is \emph{no payload}, then
\[
\begin{aligned}
& M(0) = M_r + m_p(0) \\
& M(t) = M_r + m_p(t)
\end{aligned}
\]

\emph{But}, if payload is considered, observe that payload can be considered
to be part of the whole of the vehicle (imagine an imaginary box surrounding
the entire system or ``vehicle''), \emph{and also} can be considered to be
separate and distinct (no overlap) from the propulsion system plus propellant
(imagine an imaginary box surrounding only the rocket (propulsion system) and
fuel, and another imaginary box surrounding the payload; the two imaginary
boxes do not overlap).

\[
\begin{aligned}
& M(0) = M_r + M_{\text{pay}} + m_p(0) \\
& M(t) = M_r + M_{\text{pay}} + m_p(t)
\end{aligned}
\]

Recall that mass ratio is defined as $\mathbf{MR} := \frac{M_f}{ M(0)}$. But
remember that $\mathbf{MR}$ can be considered for different, discrete
imaginary boxes (subsystems). So for Example 2-1 of Biblarz and Sutton:

\[
\begin{aligned}
&	\mathbf{MR}_{\text{tot}} = \frac{M_{f, \text{tot}}}{M(0)}
 = \frac{M_r + M_{\text{pay}}}{ M(0) } \\
	& \mathbf{MR}_r = \frac{ M_r }{ M_r+ m_p(0)}
 = \frac{M_r }{ M(0) - M_{\text{pay}}}
\end{aligned}
\]
Notice how you can not take the payload mass into account; imagine only
dealing with what's in the imaginary box surrounding only the rocket plus
propellant.

If $\frac{\partial \mathbf{u}}{ \partial t} = 0$, assuming flat space,
$F = \int_S \dot{m}_S \otimes \mathbf{u}$. Notice that for the units to make
sense, $\dot{m}_S$ is the mass per unit (surface) area. If
$\dot{m}_S , \mathbf{u}$ constant over surface area $S$,
$F= \dot{m} \mathbf{u}$.  Then, $u = \frac{F}{\dot{m}} = I_s g_0$.

Then $\overline{u} = I_s g_0$. \\
$I_{\text{tot}} = I_s m_p(0) = I_s(M(0) - (M_r + M_{\text{pay}}))$

Given the burn duration $T$, estimate $\dot{m}_{p, \text{exp}}$ with the
simplest (linear) approximation:
$\dot{m}_{p, \text{exp}} = \frac{m_p(0)}{T}$.

$\overline{F}_{\text{thrust}} = \overline{ \dot{m}}_{p, \text{exp}} g_0 I_s$

Consider $\overline{F}_{\text{thrust}}$ on rocket, only. For a horizontal
trajectory, max. acceleration is found at end of the thrusting schedule, just
before shutdown (because while thrust is unchanged the mass is now at its
minimum value). Think of ``Max-Q''.

\[
\Longrightarrow a_f = \frac{ \overline{F}_{\text{thrust}} }{ M_f}
\]

\paragraph*{Annotation: the example with numbers, in the box notation.}
The ninth edition's data (checked against the source text layer) are
$M(0)=200\ \mathrm{kg}$, $M_{f,\text{tot}}=130\ \mathrm{kg}$,
$M_{\text{pay}}=110\ \mathrm{kg}$, $T=3.0\ \mathrm{s}$, $I_s=240\ \mathrm{s}$,
and Sutton rounds $g_0$ to $9.81\ \mathrm{m/s^2}$ throughout the example.
Then $m_p(0)=200-130=70\ \mathrm{kg}$ and $M_r=130-110=20\ \mathrm{kg}$, and
the two boxes give
\[
\begin{aligned}
 \mathbf{MR}_{\text{tot}}&=\frac{130}{200}=0.65, &
 \zeta_{\text{tot}}&=\frac{70}{200}=0.35, \\
 \mathbf{MR}_r&=\frac{20}{90}=0.222, &
 \zeta_r&=\frac{70}{90}=0.778,
\end{aligned}
\]
each pair summing to one as Eq.~\ref{Eq:ZetaPlusMR} requires. Sutton
reports $\mathbf{MR}$ for both boxes but $\zeta$ only for the propulsion
unit, although the problem statement asks for both. Continuing,
$c=I_sg_0=240\times9.81=2354.4\ \mathrm{m/s}$
(with the defined $g_0=9.80665$ it would be $2353.6$; the book's rounding is
harmless here), $\overline{\dot m}=70/3\ \mathrm{kg/s}$, and
$\overline F=\overline{\dot m}\,c=(70/3)\times2354.4=54{,}936\ \mathrm{N}$
--- using the unrounded flow rate, which is what reproduces the printed
figure; the book's own display of this line is the erratum recorded in
Section~\ref{topic:force-one-form}. Then $I_t=\overline F\,T=164{,}808\
\mathrm{N\,s}$, equal to $I_sm_p(0)g_0=240\times70\times9.81$ as it must be,
and the impulse-to-weight ratio of the propulsion unit is
\[
 \frac{I_t}{w_0}=\frac{164{,}808}{90\times9.81}=186.7\ \mathrm{s}
 =I_s\,\zeta_r=240\times0.778 ,
\]
confirming Eq.~\ref{Eq:ImpulseToWeightIsIsZeta} and Sutton's printed
$187\ \mathrm{s}$. For the acceleration the box is no longer a choice: the
$35\,g_0$ limit is on the payload's electronics, which ride with the whole
vehicle, so $M_f$ in the manuscript's last line must be the vehicle's
$130\ \mathrm{kg}$ and not the propulsion unit's $20\ \mathrm{kg}$. That gives
$a_f=54{,}936/130=422.6\ \mathrm{m/s^2}=43.08\,g_0$, over the limit, and
Eq.~\ref{Eq:BurnoutAcceleration} checks it from the other end:
$a_0=54{,}936/200=274.7\ \mathrm{m/s^2}=28.0\,g_0$ and
$a_0/\mathbf{MR}_{\text{tot}}=28.0/0.65=43.1\,g_0$. The rule the example
teaches is therefore: for ratios, choose the box and say which; for
Newton's law, the box is whatever moves rigidly together.

\paragraph*{Annotation: the linear estimate is exact for the average.}
``The simplest (linear) approximation'' $\dot m_{p,\text{exp}}=m_p(0)/T$
understates its own status. By the fundamental theorem of calculus,
$\frac1T\int_0^T\dot m_{p,\text{exp}}\,dt=m_{p,\text{exp}}(T)/T=m_p(0)/T$
whatever the flow history, so as a statement about the \emph{average} flow
rate it is exact. What is an approximation is the next step: using the
average in place of the instantaneous $\dot m(T^-)$ to obtain the thrust at
burnout, which is the constant-thrust assumption A3 again. A motor with a
regressive burn has its lowest thrust at the end, and the true burnout
acceleration is lower than $\overline F/M_f$; a progressive burn, higher.

\paragraph*{Annotation: not Max-Q.}
The last sentence before the display should not invoke Max-Q. Maximum dynamic
pressure, $q=\tfrac12\rho v^2$, occurs where the product of atmospheric
density and flight speed peaks, typically early in an ascent while the air is
still dense, and it is a structural-load event, not a peak in acceleration.
Peak acceleration under constant thrust occurs at burnout because the mass is
least, which is what the sentence correctly says. The two events are
unrelated in general, and in Example~2--1 there is no atmosphere modelled at
all --- gravity is balanced by lift and drag is neglected --- so the example
has no Max-Q to think of.

\subsection{The card: Section 2.1 on one page}
Every formula of Sutton's Section~2.1, with the hypothesis under which it is
exact, so that the shortcut and its price sit side by side.

\begin{center}\small
\renewcommand{\arraystretch}{1.35}
\begin{tabular}{@{}p{0.19\linewidth}p{0.07\linewidth}p{0.30\linewidth}p{0.36\linewidth}@{}}
\hline
Impulse and velocity & Sutton & Formula & Exact when \\
\hline
Total impulse & 2--1 & $I_t=\int_0^{t}F\,dt$ &
 Definition. Equals $\lVert\Delta p\rVert$ under A1--A2
 (Props.~\ref{prop:impulse-momentum}, \ref{prop:scalar-bound}); a bound
 otherwise. \\
 & 2--2 & $I_t=Ft$ & $F$ constant, transients negligible (A3--A4). \\
Specific impulse & 2--3 & $I_s=\dfrac{\int F\,dt}{g_0\int\dot m\,dt}$ &
 Definition; the mass-flow-weighted mean of $F/(\dot mg_0)$
 (Prop.~\ref{prop:weighted-mean}). \\
 & 2--4 & $I_s=I_t/(m_pg_0)$ & Identity from 2--3. \\
 & 2--5 & $I_s=F/(\dot mg_0)$ & $F$, $\dot m$ constant. \\
Effective exhaust velocity & 2--6 & $c=I_sg_0=F/\dot m$ &
 Definition of $c$; equals $v_2$ iff $p_2=p_3$ (Eq.~2--15). \\
\hline
\end{tabular}
\end{center}

\begin{center}\small
\renewcommand{\arraystretch}{1.35}
\begin{tabular}{@{}p{0.19\linewidth}p{0.07\linewidth}p{0.30\linewidth}p{0.36\linewidth}@{}}
\hline
Masses and ratios & Sutton & Formula & Exact when \\
\hline
Mass ratio & 2--7 & $\mathbf{MR}=m_f/m_0$ & Definition; state the box. \\
Propellant mass fraction & 2--8, 2--9 & $\zeta=m_p/m_0=1-\mathbf{MR}$ &
 Same box as $\mathbf{MR}$; residual convention stated. \\
Impulse-to-weight & 2--11 & $\dfrac{I_t}{w_0}=\dfrac{I_s}{m_f/m_p+1}=I_s\zeta$ &
 Identity given 2--4; $\le I_s$. \\
Thrust-to-weight & --- & $\dfrac{F}{w_0}=\dfrac{a_0}{g_0}$,\quad
 $\dfrac{a_f}{g_0}=\dfrac{F/w_0}{\mathbf{MR}}$ &
 $F$ constant, gravity and drag neglected. \\
\hline
\end{tabular}
\end{center}

\subsection{Scope of what is claimed}
Everything above is definition, identity or bookkeeping, checkable by hand;
the control-volume derivation of Eq.~\ref{Eq:ThrustEquation} is exact under
the five assumptions labelled A\textsubscript{s}--A\textsubscript{g}, and the
numbers in Example~2--1 were recomputed from the ninth edition's data and
agree with its printed results. No executable study accompanies this topic.
The covector-valued-form version of the momentum balance, the object the
manuscript's $D$ points at, remains deferred to a Sutton~2.2 topic, as in
Section~\ref{topic:force-one-form}.

\subsubsection*{Transcription notes}
The manuscript is reproduced as written apart from the following, each a
typographical fix rather than a change of content: a missing ``$=$'' in
Eq.~\ref{Eq:SpecificImpulsePractical}; $g\to g_0$ in the second quick line
and in Eq.~\ref{Eq:ImpulseToWeight}; $M_s\to M_f$ in the last member of
Eq.~\ref{Eq:ImpulseToWeight}; ``$>0$'' $\to$ ``$\ge0$'' for
$\dot m_{p,\text{exp}}$, for the reason given in the regularity annotation;
straight quotation marks set as \LaTeX{} quotes; equation labels added so the
annotations can cite by number; and trailing forced line breaks removed where
they preceded a blank line.
